\begin{frame}[label=q3c3]{3c.3 · Decimales: distribución mensual e índice}
  \small
  \pregunta{¿Qué meses superan una distribución uniforme y cómo se concentra el total por trimestre?}
  \entrada{RDD; $m$ = mes, $t$ = trimestre, $c$ = cantidad. Periodo 2019--2024; referencia: $100/12=8.33\%$.}
  \begin{columns}[T,onlytextwidth]
    \column{0.55\textwidth}
    \textbf{Map y Reduce de cada operación}
    \parmr{A: mes}{$(m,c)$\\Clave: mes}
      {\texttt{reduceByKey(\_+\_)}\\$(m,T_m)$: total mensual}
    \parmr{B: nacional}{$(\texttt{TOTAL},c)$\\La misma clave para todas las filas}
      {\texttt{reduceByKey(\_+\_)}\\$T=121\,427$}
    \parmr{C: apoyo}{$(t,c)$\\Calcular trimestre desde el mes}
      {\texttt{reduceByKey(\_+\_)}\\$(t,S_t)$: total trimestral}
    \column{0.42\textwidth}
    \textbf{Salida}\\[5pt]
    \footnotesize
    \begin{tabular}{@{}lrrr@{}}
      \toprule Mes & Personas & \% & Índice\\\midrule
      Diciembre & 11\,522 & 9.49 & 1.14\\
      Febrero & 11\,118 & 9.16 & 1.10\\
      Abril & 8\,805 & 7.25 & 0.87\\\bottomrule
    \end{tabular}
    \par\vspace{6pt}\textbf{Cálculo después de reducir}\\
    Participación: $p_m=100T_m/T$.\\
    Índice: $I_m=12T_m/T$.
    \par\vspace{5pt}Apoyo C: trimestre 4, 32\,650 personas (26.89\%).
    \par\vspace{2pt}
    \codigo{rdd/03c_3_estacionalidad_mensual}{20}{20}
  \end{columns}
  \respuesta{Diciembre: índice 1.14, superior a la referencia uniforme de 1.}
  \origen{rdd}{03c_3_estacionalidad_mensual}
  \note{35 s. Leer primero el Map y después el Reduce. B reúne todo bajo TOTAL; C es el apoyo trimestral. A y B alimentan los cálculos decimales.}
\end{frame}
